% ============================================================
% CATÁLOGO EUF 2014-2 — Classificação Taxônomica (R-Codes)
% Dataset: Exames de Pós-Graduação (Mestrado)
% Gerado em: Maio 2026
% ============================================================
\documentclass[12pt,a4paper]{article}
\usepackage[utf8]{inputenc}
\usepackage[T1]{fontenc}
\usepackage{lmodern}
\usepackage[brazilian]{babel}
\usepackage{geometry}
\geometry{margin=2.5cm,top=2.5cm,bottom=2cm}
\usepackage{indentfirst}
\usepackage{setspace}
\onehalfspacing
\usepackage{amsmath,amssymb}
\usepackage{booktabs,array,longtable}
\usepackage{hyperref}
\usepackage{graphicx}
\usepackage{float}
\usepackage{xcolor}
\usepackage{caption}

\hypersetup{colorlinks=true,linkcolor=blue,citecolor=blue,urlcolor=blue}

\definecolor{azulprof}{RGB}{0,51,102}
\definecolor{verdeexame}{RGB}{0,100,50}

% Comando para destacar R-codes
\newcommand{\rcode}[1]{\textcolor{azulprof}{\textbf{#1}}}
\newcommand{\bloom}[1]{\textcolor{verdeexame}{\textbf{#1}}}

\title{\textbf{Catálogo da EUF 2014-2:}\\
\large Classificação Taxônomica de 10 Questões de Física\\
segundo a Taxonomia 350 de Raciocínios Científicos}
\author{OpenCode Ecosystem --- Análise de Dados\\
Mestrado PPGTE/CT/UFC}
\date{Maio 2026}

\begin{document}
\maketitle

\begin{abstract}
Este catálogo documenta a classificação das 10 questões do Exame Unificado
de Física (EUF) do segundo semestre de 2014 segundo a Taxonomia 350 de
Raciocínios Científicos. Cada questão é analisada quanto ao seu raciocínio
primário (dominante) e secundário (complementar), área de conhecimento,
categoria temática, nível taxonômico de Bloom e domínio conceitual.
O catálogo inclui tabelas de distribuição por área, por tipo de raciocínio,
por nível Bloom e um mapa de correspondência entre as questões e os 9
códigos de raciocínio mobilizados no exame.
\vspace{0.3cm}

\noindent\textbf{Palavras-chave}: EUF 2014-2. Taxonomia 350. Raciocínio
científico. Classificação de questões. Física.
\end{abstract}

\tableofcontents
\newpage

% ============================================================
\section{Visão Geral do Exame}
% ============================================================

O Exame Unificado de Física (EUF) do segundo semestre de 2014 é composto
por 10 questões distribuídas em 6 áreas da Física:

\begin{itemize}
  \item \textbf{Eletromagnetismo} (2 questões): Q01 (Capacitor esférico),
    Q02 (Bobinas de Helmholtz);
  \item \textbf{Física Moderna} (2 questões): Q03 (Radiação de corpo negro),
    Q04 (Colisão relativística inelástica);
  \item \textbf{Termodinâmica} (1 questão): Q05 (Equação de estado e energia molar);
  \item \textbf{Mecânica Clássica} (2 questões): Q06 (Pêndulo com suporte móvel),
    Q07 (Átomo de trítio e decaimento beta);
  \item \textbf{Mecânica Quântica} (2 questões): Q08 (Hamiltoniano da amônia),
    Q09 (Oscilador harmônico 3D);
  \item \textbf{Mecânica Estatística} (1 questão): Q10 (Gás ideal canônico).
\end{itemize}

As questões foram classificadas segundo 9 códigos distintos da Taxonomia
350, distribuídos em 4 categorias taxonômicas.

% ============================================================
\section{Classificação das Questões}
% ============================================================

A Tabela~\ref{tab:questoes} apresenta a classificação completa das 10
questões, incluindo: identificador, título, área primária e secundária,
raciocínio primário e secundário (com código e descrição), categoria
temática, subcategoria, nível de Bloom e domínio conceitual.

\small
\begin{longtable}{@{}p{1.5cm}p{2.8cm}p{2.5cm}p{2.5cm}p{5.2cm}@{}}
\caption{Classificação completa das 10 questões da EUF 2014-2}
\label{tab:questoes}\\
\toprule
\textbf{Questão} & \textbf{Título} & \textbf{R-primário} & \textbf{R-secundário} & \textbf{Domínio} \\
\midrule
\endfirsthead

\caption[]{Classificação completa das 10 questões da EUF 2014-2 (continuação)}\\
\toprule
\textbf{Questão} & \textbf{Título} & \textbf{R-primário} & \textbf{R-secundário} & \textbf{Domínio} \\
\midrule
\endhead

\bottomrule
\endfoot

Q01 & Capacitor esférico
  & \rcode{R065} (Simetria-Conservacao)
  & \rcode{R069} (Variacional-Principio)
  & Capacitor esférico: Lei de Gauss, simetria esférica,
    densidade de energia, capacitância \\
\midrule

Q02 & Bobinas de Helmholtz
  & \rcode{R065} (Simetria-Conservacao)
  & \rcode{R069} (Variacional-Principio)
  & Bobinas de Helmholtz: simetria axial, condição de campo
    uniforme, $dB/dz=0$ e $d^2B/dz^2=0$ \\
\midrule

Q03 & Radiação de corpo negro
  & \rcode{R070} (Limite-Assintotico)
  & \rcode{R071} (Dualidade-Transformacao)
  & Radiação de corpo negro: limites assintóticos
    $h\nu\ll kT$, transformação $\nu\to\lambda$, lei de
    Stefan-Boltzmann \\
\midrule

Q04 & Colisão relativística inelástica
  & \rcode{R071} (Dualidade-Transformacao)
  & \rcode{R065} (Simetria-Conservacao)
  & Colisão inelástica relativística: transformação de
    Lorentz, conservação de 4-momento, perda de energia \\
\midrule

Q05 & Equação de estado e energia molar
  & \rcode{R008} (Deducao-Formal-Passo-a-Passo)
  & \rcode{R066} (Dimensional-Analitico)
  & Equação de estado e energia molar: derivadas parciais,
    identidade $(\partial u/\partial v)_T = T(\partial p/\partial T)_v - p$,
    expansão livre \\
\midrule

Q06 & Pêndulo com suporte móvel
  & \rcode{R069} (Variacional-Principio)
  & \rcode{R054} (Otimizacao-Restrita)
  & Pêndulo com suporte móvel: Lagrangiana, Euler-Lagrange,
    oscilações forçadas, frequência de ressonância \\
\midrule

Q07 & Átomo de trítio e decaimento beta
  & \rcode{R065} (Simetria-Conservacao)
  & \rcode{R069} (Variacional-Principio)
  & Trítio e decaimento beta: força central, conservação de
    energia/momento angular, transição de órbita \\
\midrule

Q08 & Hamiltoniano da amônia
  & \rcode{R109} (Espectral-Operadores)
  & \rcode{R008} (Deducao-Formal-Passo-a-Passo)
  & Hamiltoniano da amônia: diagonalização $2\times2$,
    autovalores, probabilidade de transição, regra de Born \\
\midrule

Q09 & Oscilador harmônico 3D
  & \rcode{R010} (Decomposicao-Modular)
  & \rcode{R109} (Espectral-Operadores)
  & Oscilador harmônico 3D: separação de variáveis,
    equação radial, espectro, degenerescência \\
\midrule

Q10 & Gás ideal canônico
  & \rcode{R069} (Variacional-Principio)
  & \rcode{R070} (Limite-Assintotico)
  & Gás ideal canônico: Hamiltoniana, integral de fase,
    função de partição, limite termodinâmico \\
\bottomrule
\end{longtable}

\newpage
\section{Mapa de Raciocínios}

\subsection{Códigos Utilizados}

A Tabela~\ref{tab:rcodes} lista os 9 códigos de raciocínio mobilizados
no exame, com suas categorias taxonômicas e descrições.

\begin{table}[H]
\centering
\caption{Códigos de raciocínio utilizados na EUF 2014-2}
\label{tab:rcodes}
\begin{tabular}{cllc}
\toprule
\textbf{Código} & \textbf{Nome} & \textbf{Categoria} & \textbf{Frequência} \\
\midrule
\rcode{R008} & Deducao-Formal-Passo-a-Passo
  & I. Lógico-Dedutivos & 2 (Q05, Q08) \\
\rcode{R010} & Decomposicao-Modular
  & II. Analítico-Estruturais & 1 (Q09) \\
\rcode{R054} & Otimizacao-Restrita
  & VII. Econômico-Decisórios & 1 (Q06) \\
\rcode{R065} & Simetria-Conservacao
  & XIII. Físico-Matemáticos & 4 (Q01, Q02, Q04, Q07) \\
\rcode{R066} & Dimensional-Analitico
  & XIII. Físico-Matemáticos & 1 (Q05) \\
\rcode{R069} & Variacional-Principio
  & XIII. Físico-Matemáticos & 5 (Q01, Q02, Q06, Q07, Q10) \\
\rcode{R070} & Limite-Assintotico
  & XIII. Físico-Matemáticos & 2 (Q03, Q10) \\
\rcode{R071} & Dualidade-Transformacao
  & XIII. Físico-Matemáticos & 2 (Q03, Q04) \\
\rcode{R109} & Espectral-Operadores
  & XIV. Lógica Matemática Avançada & 2 (Q08, Q09) \\
\bottomrule
\end{tabular}
\end{table}

\subsection{Distribuição por Categoria Taxonômica}

A Tabela~\ref{tab:categorias} mostra a distribuição das ocorrências de
raciocínio (primário + secundário) pelas categorias da Taxonomia 350.

\begin{table}[H]
\centering
\caption{Distribuição por categoria taxonômica}
\label{tab:categorias}
\begin{tabular}{lcc}
\toprule
\textbf{Categoria} & \textbf{Ocorrências} & \textbf{Percentual} \\
\midrule
Categoria XIII: Físico-Matemáticos & 14 & 70\% \\
Categoria I: Lógico-Dedutivos & 3 & 15\% \\
Categoria XIV: Lógica Matemática Avançada & 2 & 10\% \\
Categoria II: Analítico-Estruturais & 1 & 5\% \\
Categoria VII: Econômico-Decisórios & 1 & 5\% \\
\midrule
\textbf{Total} & \textbf{20} & \textbf{100\%} \\
\bottomrule
\end{tabular}
\end{table}

\subsection{Distribuição por Área da Física}

A Tabela~\ref{tab:areas} cruza as áreas da física com os raciocínios
primários, revelando padrões disciplinares.

\begin{table}[H]
\centering
\caption{Distribuição áreas $\times$ raciocínio primário}
\label{tab:areas}
\begin{tabular}{lcc}
\toprule
\textbf{Área} & \textbf{Raciocínio Primário} & \textbf{Bloom} \\
\midrule
Eletromagnetismo & R065 (Simetria-Conservacao) $\times$2 & Aplicar + Analisar \\
Física Moderna & R070 (Limite-Assintotico) + R071 (Dualidade-Transformacao) & Analisar + Aplicar \\
Termodinâmica & R008 (Deducao-Formal-Passo-a-Passo) & Analisar \\
Mecânica Clássica & R069 (Variacional-Principio) + R065 (Simetria-Conservacao) & Aplicar $\times$2 \\
Mecânica Quântica & R109 (Espectral-Operadores) + R010 (Decomposicao-Modular) & Aplicar + Analisar \\
Mecânica Estatística & R069 (Variacional-Principio) & Aplicar \\
\bottomrule
\end{tabular}
\end{table}

\subsection{Níveis de Bloom}

A Tabela~\ref{tab:bloom} apresenta a distribuição dos níveis taxonômicos
de Bloom nos 10 enunciados.

\begin{table}[H]
\centering
\caption{Distribuição dos níveis de Bloom}
\label{tab:bloom}
\begin{tabular}{lccl}
\toprule
\textbf{Nível} & \textbf{Questões} & \textbf{Total} & \textbf{Percentual} \\
\midrule
\bloom{Aplicar} & Q01, Q04, Q06, Q07, Q08, Q10 & 6 & 60\% \\
\bloom{Analisar} & Q02, Q03, Q05, Q09 & 4 & 40\% \\
\midrule
\textbf{Total} & & \textbf{10} & \textbf{100\%} \\
\bottomrule
\end{tabular}
\end{table}

\subsection{Matriz de Incidência}

A Tabela~\ref{tab:matriz} mostra a matriz de incidência Questões
$\times$ Raciocínios, onde P = raciocínio primário e S = raciocínio
secundário.

\begin{table}[H]
\centering
\caption{Matriz de incidência: Questões $\times$ Raciocínios}
\label{tab:matriz}
\begin{tabular}{lccccccccc}
\toprule
\textbf{Questão} & R008 & R010 & R054 & R065 & R066 & R069 & R070 & R071 & R109 \\
\midrule
Q01 &   &   &   & P &   & S &   &   &   \\
Q02 &   &   &   & P &   & S &   &   &   \\
Q03 &   &   &   &   &   &   & P & S &   \\
Q04 &   &   &   & S &   &   &   & P &   \\
Q05 & P &   &   &   & S &   &   &   &   \\
Q06 &   &   & S &   &   & P &   &   &   \\
Q07 &   &   &   & P &   & S &   &   &   \\
Q08 & S &   &   &   &   &   &   &   & P \\
Q09 &   & P &   &   &   &   &   &   & S \\
Q10 &   &   &   &   &   & P & S &   &   \\
\midrule
\textbf{Total} & 2 & 1 & 1 & 4 & 1 & 5 & 2 & 2 & 2 \\
\bottomrule
\end{tabular}
\end{table}

% ============================================================
\section{Análise por Questão}
% ============================================================

\subsection{Eletromagnetismo}

\subsubsection{Q01 --- Capacitor Esférico (\bloom{Aplicar})}

\textbf{Enunciado:} Um capacitor esférico de raios $a$ (interno) e $b$ (externo)
é carregado com carga $Q$. Determine a capacitância, a densidade de energia
do campo elétrico e a energia total armazenada.

\medskip
\noindent\textbf{Classificação:}
\begin{itemize}
  \item \textbf{Raciocínio primário:} \rcode{R065} (Simetria-Conservacao) ---
    Aplica a Lei de Gauss em coordenadas esféricas para determinar o campo
    elétrico. A simetria esférica do problema reduz a integral de superfície
    a $E(r)\cdot 4\pi r^2 = Q_{\text{int}}/\varepsilon_0$.
  \item \textbf{Raciocínio secundário:} \rcode{R069} (Variacional-Principio) ---
    Calcula a energia total como integral da densidade de energia
    $u = \varepsilon_0 E^2/2$ sobre o volume entre as placas, tratando a
    energia como funcional do campo.
\end{itemize}

\subsubsection{Q02 --- Bobinas de Helmholtz (\bloom{Analisar})}

\textbf{Enunciado:} Duas bobinas idênticas de raio $R$ com $N$ espiras,
separadas por distância $d$, são percorridas por corrente $I$ no mesmo
sentido. Determine a condição para que o campo magnético seja máximo
uniforme na região central.

\medskip
\noindent\textbf{Classificação:}
\begin{itemize}
  \item \textbf{Raciocínio primário:} \rcode{R065} (Simetria-Conservacao) ---
    Utiliza a simetria axial do problema para aplicar a Lei de Biot-Savart
    ao longo do eixo $z$.
  \item \textbf{Raciocínio secundário:} \rcode{R069} (Variacional-Principio) ---
    Impõe as condições $dB/dz = 0$ e $d^2B/dz^2 = 0$ no ponto médio como
    condição de campo uniforme (extremal).
\end{itemize}

\subsection{Física Moderna}

\subsubsection{Q03 --- Radiação de Corpo Negro (\bloom{Analisar})}

\textbf{Enunciado:} A densidade espectral de energia da radiação de corpo
negro é dada pela Lei de Planck. Determine a constante de Stefan-Boltzmann
a partir da integração da distribuição de Planck e discuta os limites
clássicos.

\medskip
\noindent\textbf{Classificação:}
\begin{itemize}
  \item \textbf{Raciocínio primário:} \rcode{R070} (Limite-Assintotico) ---
    Expande a distribuição de Planck para $h\nu/kT \ll 1$, recuperando a
    Lei de Rayleigh-Jeans, e explora o limite de Wien para $h\nu/kT \gg 1$.
  \item \textbf{Raciocínio secundário:} \rcode{R071} (Dualidade-Transformacao) ---
    Realiza a mudança de variável $\nu \to \lambda$ para expressar a
    densidade espectral em termos de comprimento de onda e integra para
    obter $u = aT^4$.
\end{itemize}

\subsubsection{Q04 --- Colisão Relativística Inelástica (\bloom{Aplicar})}

\textbf{Enunciado:} Uma partícula de massa $m$ e energia $E$ colide
inelasticamente com uma partícula idêntica em repouso. Determine a
massa do sistema resultante e a energia perdida.

\medskip
\noindent\textbf{Classificação:}
\begin{itemize}
  \item \textbf{Raciocínio primário:} \rcode{R071} (Dualidade-Transformacao) ---
    Aplica transformações de Lorentz entre referenciais para analisar a
    colisão em diferentes sistemas de coordenadas.
  \item \textbf{Raciocínio secundário:} \rcode{R065} (Simetria-Conservacao) ---
    Utiliza a conservação do 4-momento total $P^\mu$ como expressão da
    simetria de Poincaré (invariância por translações espaço-temporais).
\end{itemize}

\subsection{Termodinâmica}

\subsubsection{Q05 --- Equação de Estado e Energia Molar (\bloom{Analisar})}

\textbf{Enunciado:} Para um gás real descrito por $p(v,T)$, deduza a
relação entre a energia interna molar e as derivadas parciais da
equação de estado. Aplique à expansão livre.

\medskip
\noindent\textbf{Classificação:}
\begin{itemize}
  \item \textbf{Raciocínio primário:} \rcode{R008} (Deducao-Formal-Passo-a-Passo) ---
    Constrói cadeia dedutiva de identidades termodinâmicas partindo da
    primeira lei e das relações de Maxwell: $(\partial u/\partial v)_T
    = T(\partial p/\partial T)_v - p$.
  \item \textbf{Raciocínio secundário:} \rcode{R066} (Dimensional-Analitico) ---
    Utiliza análise dimensional e relações entre variáveis termodinâmicas
    ($p, v, T, u, s$) para verificar consistência das expressões.
\end{itemize}

\subsection{Mecânica Clássica}

\subsubsection{Q06 --- Pêndulo com Suporte Móvel (\bloom{Aplicar})}

\textbf{Enunciado:} Um pêndulo simples de comprimento $\ell$ e massa $m$
tem seu ponto de suspensão movendo-se horizontalmente com aceleração
$a(t)$. Determine as equações de movimento e as frequências de oscilação.

\medskip
\noindent\textbf{Classificação:}
\begin{itemize}
  \item \textbf{Raciocínio primário:} \rcode{R069} (Variacional-Principio) ---
    Constrói a Lagrangiana do sistema e aplica as equações de
    Euler-Lagrange para obter as equações de movimento.
  \item \textbf{Raciocínio secundário:} \rcode{R054} (Otimizacao-Restrita) ---
    Determina a frequência de ressonância do oscilador forçado como
    condição de máxima amplitude.
\end{itemize}

\subsubsection{Q07 --- Átomo de Trítio e Decaimento Beta (\bloom{Aplicar})}

\textbf{Enunciado:} Um átomo de trítio ($^3$H) sofre decaimento beta,
transformando-se em $^3$He$^+$. Determine a energia e o momento angular
do elétron e analise a transição orbital.

\medskip
\noindent\textbf{Classificação:}
\begin{itemize}
  \item \textbf{Raciocínio primário:} \rcode{R065} (Simetria-Conservacao) ---
    Utiliza a conservação de energia e momento angular no campo de força
    central coulombiano para analisar a órbita eletrônica.
  \item \textbf{Raciocínio secundário:} \rcode{R069} (Variacional-Principio) ---
    Calcula a energia mecânica total e a trajetória via potencial efetivo.
\end{itemize}

\subsection{Mecânica Quântica}

\subsubsection{Q08 --- Hamiltoniano da Amônia (\bloom{Aplicar})}

\textbf{Enunciado:} A molécula de $NH_3$ é descrita por um Hamiltoniano
$2\times2$ com estados $|L\rangle$ e $|R\rangle$. Determine os autovalores,
autovetores e a probabilidade de transição entre os estados.

\medskip
\noindent\textbf{Classificação:}
\begin{itemize}
  \item \textbf{Raciocínio primário:} \rcode{R109} (Espectral-Operadores) ---
    Diagonaliza o operador Hamiltoniano $2\times2$ para obter autovalores
    e autovetores, fundamento da mecânica quântica.
  \item \textbf{Raciocínio secundário:} \rcode{R008} (Deducao-Formal-Passo-a-Passo) ---
    Aplica sequencialmente os postulados da MQ: operador $\to$ autovalores
    $\to$ probabilidade (regra de Born) $\to$ evolução temporal.
\end{itemize}

\subsubsection{Q09 --- Oscilador Harmônico 3D (\bloom{Analisar})}

\textbf{Enunciado:} Uma partícula de massa $m$ está sujeita ao potencial
$V(\mathbf{r}) = \frac{1}{2}m\omega^2 r^2$. Determine os autovalores de
energia, a degenerescência e a equação radial.

\medskip
\noindent\textbf{Classificação:}
\begin{itemize}
  \item \textbf{Raciocínio primário:} \rcode{R010} (Decomposicao-Modular) ---
    Separa a equação de Schrödinger em coordenadas cartesianas (três
    osciladores 1D independentes) e, alternativamente, em coordenadas
    esféricas (parte radial + harmônicos esféricos).
  \item \textbf{Raciocínio secundário:} \rcode{R109} (Espectral-Operadores) ---
    Determina o espectro de energia $E = \hbar\omega(n+3/2)$ com
    degenerescência $g(n) = (n+1)(n+2)/2$.
\end{itemize}

\subsection{Mecânica Estatística}

\subsubsection{Q10 --- Gás Ideal Canônico (\bloom{Aplicar})}

\textbf{Enunciado:} Considere $N$ partículas de gás ideal clássico em
contato com um reservatório térmico à temperatura $T$. Calcule a função
de partição canônica e as propriedades termodinâmicas.

\medskip
\noindent\textbf{Classificação:}
\begin{itemize}
  \item \textbf{Raciocínio primário:} \rcode{R069} (Variacional-Principio) ---
    Formula o problema via Hamiltoniana clássica e calcula a integral de
    fase $Z = \int e^{-\beta H} d^{3N}q\,d^{3N}p$ no ensemble canônico.
  \item \textbf{Raciocínio secundário:} \rcode{R070} (Limite-Assintotico) ---
    Aplica o limite termodinâmico $N\to\infty$, $V\to\infty$, com
    densidade $N/V$ constante, para obter propriedades extensivas.
\end{itemize}

% ============================================================
\section{Distribuição Comparativa}
% ============================================================

\subsection{Comparação com os 6 Exames EUF do Dataset}

A Tabela~\ref{tab:comparativo} compara a EUF 2014-2 com os demais 5
exames EUF presentes no dataset (todas as 10 questões por exame, nas
mesmas 6 áreas da física). Os dados revelam variações significativas
no repertório de raciocínios mobilizados.

\begin{table}[H]
\centering
\caption{Comparativo entre os 6 exames EUF do dataset}
\label{tab:comparativo}
\begin{tabular}{lcccccc}
\toprule
\textbf{Métrica} & \textbf{2011-1} & \textbf{2012-2} & \textbf{2014-2} & \textbf{2015-1} & \textbf{2015-2} & \textbf{2016-1} \\
\midrule
R-codes distintos & 8 & 7 & 7 & 5 & 7 & 8 \\
Cat. XIII (Físico-Matemáticos) & 4 & 5 & \textbf{7} & 4 & 4 & 4 \\
Áreas & 6 & 6 & 6 & 6 & 6 & 6 \\
R-codes da Cat. XIII & 4 & 5 & \textbf{7} & 4 & 4 & 4 \\
\bottomrule
\end{tabular}
\end{table}

\noindent A EUF 2014-2 destaca-se com \textbf{7 ocorrências} de
raciocínios Físico-Matemáticos (Cat. XIII), o maior valor entre todos
os exames. A média dos demais 5 exames é de 4,2 ocorrências por exame,
indicando que a EUF 2014-2 demandou substancialmente mais raciocínio
quantitativo-matemático que a média histórica.

\subsection{Comparação de Raciocínios Primários}

A Tabela~\ref{tab:comparativo_rcodes} mostra todos os R-codes utilizados
como raciocínio primário em cada exame, revelando o repertório
taxonômico de cada prova.

\small
\begin{longtable}{lcccccc}
\caption{R-codes primários por exame EUF}
\label{tab:comparativo_rcodes}\\
\toprule
\textbf{R-código} & \textbf{2011-1} & \textbf{2012-2} & \textbf{2014-2} & \textbf{2015-1} & \textbf{2015-2} & \textbf{2016-1} \\
\midrule
\endfirsthead

\caption[]{R-codes primários por exame EUF (continuação)}\\
\toprule
\textbf{R-código} & \textbf{2011-1} & \textbf{2012-2} & \textbf{2014-2} & \textbf{2015-1} & \textbf{2015-2} & \textbf{2016-1} \\
\midrule
\endhead

\bottomrule
\endfoot

R008 (Deducao-Formal-Passo-a-Passo) & & & 1 & & & \\
R010 (Decomposicao-Modular) & & & 1 & & & \\
R011 (Inducao-Estrutural) & & & & & 1 & \\
R065 (Simetria-Conservacao) & 1 & 2 & 3 & & 1 & 1 \\
R067 (Perturbativo) & & 1 & & & & \\
R069 (Variacional-Principio) & 1 & 1 & 2 & 2 & 1 & 1 \\
R070 (Limite-Assintotico) & & & 1 & & & \\
R071 (Dualidade-Transformacao) & 2 & 1 & 1 & 2 & 2 & 2 \\
R073 (Caos-Deterministico) & 1 & & & & & \\
R109 (Espectral-Operadores) & & & 1 & & & \\
R117 (Monte-Carlo-MCMC) & 1 & 2 & & 2 & 2 & 2 \\
R183 (Mecanica-Quantica-Particulas) & 3 & 2 & & 3 & 2 & 1 \\
R201 (Superposicao-Quantica) & & & & & & 1 \\
R202 (Emaranhamento-Quantico) & & & & & 1 & 1 \\
R229 (Homotopia-Fundamental) & 1 & 1 & & 1 & & 1 \\
\midrule
\textbf{Total} & \textbf{10} & \textbf{10} & \textbf{10} & \textbf{10} & \textbf{10} & \textbf{10} \\
\end{longtable}

\normalsize
\noindent Observa-se que:
\begin{itemize}
  \item \rcode{R071} (Dualidade-Transformacao) é o único R-código presente
    em todos os 6 exames;
  \item \rcode{R069} (Variacional-Principio) e \rcode{R065}
    (Simetria-Conservacao) aparecem em 5 dos 6 exames, confirmando seu
    papel central na física teórica;
  \item A EUF 2014-2 é a única que mobiliza \rcode{R008}
    (Deducao-Formal), \rcode{R010} (Decomposicao-Modular),
    \rcode{R070} (Limite-Assintotico) e \rcode{R109}
    (Espectral-Operadores);
  \item Os exames 2015-1, 2015-2 e 2016-1 incorporam R-codes de
    Mecânica Quântica Aplicada (R183, R201, R202), ausentes na
    EUF 2014-2.
\end{itemize}

\subsection{Padrões Emergentes}

A análise das 10 questões revela três padrões importantes:

\begin{enumerate}
  \item \textbf{Predomínio de Raciocínios Físico-Matemáticos (Cat. XIII):}
    70\% das ocorrências pertencem a esta categoria, refletindo a natureza
    essencialmente quantitativa e matemática da física universitária.
    Destacam-se \rcode{R069} (Variacional-Principio) com 5 ocorrências
    e \rcode{R065} (Simetria-Conservacao) com 4.

  \item \textbf{Complementaridade R065-R069:} Em 4 questões (Q01, Q02, Q07),
    os raciocínios de Simetria-Conservacao e Variacional-Principio
    aparecem como par primário-secundário. Esta dupla representa o núcleo
    metodológico da física: identificar simetrias $\to$ formular
    funcional/variacional $\to$ extrair equações.

  \item \textbf{Bloom concentrado em Aplicar (60\%):} A predominância do
    nível ``Aplicar'' indica que a EUF 2014-2 privilegia a aplicação de
    métodos consagrados a problemas específicos, em vez de demandar
    análise crítica ou síntese criativa.
\end{enumerate}

% ============================================================
\section{Conclusão do Catálogo}
% ============================================================

O catálogo aqui apresentado documenta a classificação taxonômica completa
da EUF 2014-2, estabelecendo uma ponte entre o conteúdo programático do
exame e a Taxonomia 350 de Raciocínios Científicos. As 10 questões
mobilizam 9 tipos de raciocínio em 5 categorias, com forte concentração
na Categoria XIII (Físico-Matemáticos).

Esta classificação alimenta o dataset de 150 questões do projeto de
mestrado, permitindo análises estatísticas, correlações entre áreas e
raciocínios, e a validação do modelo taxonômico em domínios da física
teórica.

\medskip
\noindent\textbf{Nota:} A análise completa dos 6 exames EUF e 3 exames ENA
(totalizando 150 questões) encontra-se no relatório
\texttt{analise\_mapa\_taxonomico.pdf}, que inclui:
\begin{itemize}
  \item Correlações $\chi^2$ e V de Cramer entre área, R-code e Bloom;
  \item Comparação detalhada EUF 2014-2 vs. demais exames;
  \item Distribuição de 32 R-codes em 5 categorias taxonômicas;
  \item Classificação Bloom completa das 150 questões.
\end{itemize}

% ============================================================
% APÊNDICE: GLOSSÁRIO DE R-CODES
% ============================================================
\newpage
\appendix
\section{Glossário de R-Codes da Taxonomia 350}

\small
\begin{longtable}{@{}p{2.2cm}p{3.5cm}p{2.5cm}p{5.5cm}@{}}
\caption{Glossário dos 9 R-codes utilizados na EUF 2014-2}
\label{tab:glossario}\\
\toprule
\textbf{Código} & \textbf{Nome} & \textbf{Categoria} & \textbf{Descrição} \\
\midrule
\endfirsthead

\caption[]{Glossário (continuação)}\\
\toprule
\textbf{Código} & \textbf{Nome} & \textbf{Categoria} & \textbf{Descrição} \\
\midrule
\endhead

\bottomrule
\endfoot

\rcode{R008} & Deducao-Formal-Passo-a-Passo
  & I. Lógico-Dedutivos
  & Encadeamento lógico-dedutivo onde cada passo segue necessariamente
    do anterior por regras de inferência explícitas. \\
\rcode{R010} & Decomposicao-Modular
  & II. Analítico-Estruturais
  & Fragmentação de um sistema complexo em subsistemas independentes
    que podem ser analisados separadamente. \\
\rcode{R054} & Otimizacao-Restrita
  & VII. Econômico-Decisórios
  & Busca por extremos (máximo/mínimo) de uma função sujeita a
    vínculos ou restrições. \\
\rcode{R065} & Simetria-Conservacao
  & XIII. Físico-Matemáticos
  & Identificação de simetrias no sistema para derivar leis de
    conservação (teorema de Noether) e simplificar equações. \\
\rcode{R066} & Dimensional-Analitico
  & XIII. Físico-Matemáticos
  & Análise de relações entre grandezas via suas dimensões físicas
    fundamentais (massa, comprimento, tempo, etc.). \\
\rcode{R069} & Variacional-Principio
  & XIII. Físico-Matemáticos
  & Formulação de problemas físicos via princípios de extremo (mínima
    ação, Hamilton, Euler-Lagrange) e funcionais. \\
\rcode{R070} & Limite-Assintotico
  & XIII. Físico-Matemáticos
  & Análise do comportamento de funções ou sistemas em regimes
    extremos ($x\to 0$, $x\to\infty$, $N\to\infty$, etc.). \\
\rcode{R071} & Dualidade-Transformacao
  & XIII. Físico-Matemáticos
  & Transformação entre representações matemáticas ou referenciais
    (Lorentz, Fourier, coordenadas) para simplificar problemas. \\
\rcode{R109} & Espectral-Operadores
  & XIV. Lógica Matemática Avançada
  & Determinação de autovalores e autovetores de operadores lineares
    para caracterizar sistemas dinâmicos. \\
\bottomrule
\end{longtable}

\end{document}
